% Capstone 1: big-integer arithmetic. Port of «Мини-курсовая - арифметика».
% The performance numbers come from the generated requirements.tex (make bench).
\documentclass{ca-capstone}
\capstonerefs{gmp,gmpy2,crandall-pomerance,welschenbach,knuth-taocp2}
\begin{document}
% Generated by `make bench CAPSTONE=bigint-arithmetic` (taskgen.capstone report) from
% bench/results.json: reference time x slack on the reference machine.
% Do not edit: change bench/config.yaml and re-run the benchmark.
\newcommand{\ReqSlack}{10}
\newcommand{\ReqBits}{2048}
\newcommand{\ReqMulUs}{50}
\newcommand{\ReqDivUs}{50}
\newcommand{\ReqGcdMs}{10}
\newcommand{\ReqPowmodMs}{100}

\capstonetitle{Арифметика больших чисел}

\goal
Требуется реализовать библиотеку для работы с большими числами, позволяющую выполнять базовые
арифметические операции над знаковыми целыми числами
\emph{произвольной}\footnote{можно считать, что битовая длина любого числа ограничена сверху
некоторым значением, например, $512 \cdot 4 \cdot 8 = 16384$~бит (\emph{как получено это
значение?})} длины.

\functionality
\begin{enumerate}
  \item Сложение и вычитание
  \item Умножение
    \begin{enumerate}[label=\alph*.]
      \item школьный метод "--- <<в столбик>>
      \item метод Карацубы
    \end{enumerate}
  \item Деление
    \begin{enumerate}[label=\alph*.]
      \item нацело с остатком: функция должна возвращать частное и остаток
      \item взятие остатка от деления
    \end{enumerate}
  \item Алгоритм Евклида и расширенный алгоритм Евклида\footnote{необходимы для нахождения
    значений НОД и мультипликативно обратных, соответственно}
  \item Операции сдвига (деление/умножение на степени~2)
  \item Редукция (преобразование) Монтгомери и модульное умножение методом Монтгомери
  \item Быстрые алгоритмы возведения в степень
    \begin{enumerate}[label=\alph*.]
      \item \emph{бинарный} алгоритм: схемой <<слева-направо>>\footnote{схема
        <<слева-направо>> вычислительно более эффективная (\emph{почему?})}
      \item $q$-арный алгоритм возведения в степень: показатель раскладывается по
        основанию~$q$, принимающему одно из заранее определённых
        значений\footnote{значения $q$ должны быть степенью~2 (\emph{почему?})} (например,
        $\{2, 4, 8, 16\}$)
    \end{enumerate}
  \item Быстрое модульное возведение в степень: бинарный и $q$-арный алгоритмы для возведения
    числа в степень по модулю, использующие умножение методом Монтгомери
\end{enumerate}

\requirements
Очевидными требованиями являются:
\begin{enumerate}
  \item корректность результата выполняемых операций
  \item достаточно высокая скорость работы (см. комментарий ниже)
  \item принципиальная возможность работы с числами произвольной длины
\end{enumerate}

\textbf{Корректность работы} можно проверять, сравнивая результаты работы своих функций с
результатами функций некоторой библиотеки, например, GNU GMP~\cite{gmp}.

Приемлемой будет считаться \textbf{скорость работы}, позволяющая выполнять базовые
арифметические операции с числами длиной примерно
\textbf{\ReqBits~бит}\footnote{\emph{2048} "--- текущая стандартная длина чисел,
использующихся в RSA}. Ориентиры для чисел такой длины (одно ядро современного процессора
\emph{x86-64}):
\begin{itemize}
  \item умножение двух \ReqBits-битных чисел "--- не более \textbf{\ReqMulUs~мкс};
  \item деление с остатком вдвое более длинного числа на \ReqBits-битное "--- не более
    \textbf{\ReqDivUs~мкс};
  \item НОД, расширенный алгоритм Евклида, мультипликативно обратный "--- не более
    \textbf{\ReqGcdMs~мс};
  \item модульное возведение в степень (\ReqBits-битные модуль и показатель) "--- не более
    \textbf{\ReqPowmodMs~мс}.
\end{itemize}

\details
В работе требуется реализовать базовые операции <<с нуля>>, поэтому, очевидно, запрещено
пользоваться соответствующими функциями специализированных библиотек.

Для хранения больших чисел потребуется использовать структуры или классы, <<цифры>> больших
чисел следует хранить в массиве переменных целочисленного типа.

Так как реализация операций должна быть относительно эффективной, желательно использовать
компилируемые языки вроде \emph{C}, \emph{C++}, \emph{Go}, \emph{Swift}, \emph{Rust} и~т.\,п.

Впрочем, \textbf{язык может быть любым}. Например, если ваш класс на \emph{php} или
\emph{JS}\footnote{в этом месте автор перекрестился} достигает требуемого уровня
производительности, выбор языка будет обоснован.

Платформо-независимый, \textbf{переносимый код} писать \textbf{не требуется}, но во время
сдачи должна быть возможность проверить работу библиотеки, так что желательно
ориентироваться на \textbf{\emph{x86-64}} без расширения \emph{AVX}.

\schedule
Работа выполняется \textbf{индивидуально}. Срок выполнения "--- \textbf{месяц}. Работа должна
быть ваша: реализация должна быть \textbf{разработана и написана вами}, а не собрана из кусков
чужих работ.

Желательно хранить код в \emph{git}-репозитории на \emph{github}, \emph{bitbucket} или подобных
платформах. \textbf{Документация не требуется}.

\remarks
\begin{enumerate}
  \item Возведение числа в квадрат можно выполнять быстрее обычного умножения
  \item На языке ассемблера сдвиги и переносы в соседний значащий бит (который может
    располагаться в следующем элементе массива цифр) выполнять гораздо удобнее и
    эффективнее\footnote{какая инструкция выполняет сложение с учётом флага переноса?}
  \item Операции <<по Монтгомери>> требуют предварительной настройки (вычисления некоторых
    величин), которую достаточно выполнить один раз. Следовательно, необходимую информацию
    можно хранить в некотором хранилище (структуре, экземпляре класса,~\ldots)
  \item Проверку корректности работы реализуемых операций можно осуществлять с помощью т.\,н.
    \emph{Unit-тестов}
\end{enumerate}

\references
\end{document}
